\chapter{Optimisation case studies}
\label{ch:optcases}

\begin{aside}
Three real cases where \maya{} apps had a performance
problem, and how each was solved. Specific numbers from
specific bug reports.
\end{aside}

\section{Case 1: typing lag in a 5000-line buffer}

\textbf{Symptom.} Each keystroke had a 50 ms latency.

\textbf{Diagnosis.} The whole buffer was re-rendered every
keystroke. Layout walked 5000 elements.

\textbf{Fix.} Switched to virtualised rendering. Only the
visible 50 rows render. Latency dropped to under 2 ms.

\section{Case 2: scrolling stutter}

\textbf{Symptom.} Scrolling a list felt jerky.

\textbf{Diagnosis.} Each row computed an expensive derived
value on every render. The derivation depended on a list
that was stable.

\textbf{Fix.} Hoisted the derivation to a memo at parent
scope. Computed once, reused per row. Scrolling became
smooth.

\section{Case 3: 100\% CPU on idle}

\textbf{Symptom.} The app pegged a core even when not in
use.

\textbf{Diagnosis.} A signal updated every frame
(clock-based). An effect that wrote a derived signal read
this clock. The chain caused infinite re-renders.

\textbf{Fix.} Removed the unnecessary clock dependency from
the effect. Idle CPU returned to zero.

\section{Common themes}

\begin{itemize}[leftmargin=*]
  \item Virtualisation for big lists.
  \item Hoist expensive memos.
  \item Watch for accidental clock dependencies.
\end{itemize}

\section{Recap}

\begin{recap}
\begin{itemize}[leftmargin=*]
  \item Real cases, real numbers.
  \item Virtualisation, memo hoisting, dependency hygiene
    are the big three.
\end{itemize}
\end{recap}

\section*{Workshop}

\begin{enumerate}[leftmargin=*]
  \item Find a slow path in your app; apply one of the
    three techniques.
  \item Measure before and after.
  \item Add a regression test.
\end{enumerate}
